\chapter{Generative Media: Gambar, Video, \& Audio}
\label{chap:gen_media}

AI tidak lagi bisu dan buta. Generasi multimodal adalah tren masa depan.

\section{Diffusion Models (Stable Diffusion)}

Bagaimana AI menggambar?
Prosesnya disebut **Denoising**. AI dilatih dengan mengambil gambar bersih, lalu menambahkan "noise" (bintik-bintik buram) sedikit demi sedikit sampai menjadi acak total.
Kemudian, AI belajar membalikkan prosesnya: dari bintik acak menjadi gambar jelas.

\section{Video Generation (Sora, Runway)}

Video adalah sekumpulan gambar yang bergerak. Tantangannya adalah **Konsistensi Temporal**.
Jika AI membuat frame 1 (kucing duduk) dan frame 2 (kucing berdiri) tanpa transisi mulus, video akan terlihat berkedip (flickering). Model terbaru seperti Sora mengatasi ini dengan memahami fisika dunia nyata secara 3D.

\section{Audio \Audio & Voice Cloning Voice Cloning}

Teknologi TTS (Text-to-Speech) kini sudah sulit dibedakan dari manusia.

\begin{codeblock}[title=Bahaya Voice Cloning]
Hanya dengan sampel suara 3 detik, AI bisa meniru suara siapa pun. Ini memicu gelombang penipuan baru via telepon.
\end{codeblock}

\section{Tutorial: Menggunakan OpenAI Whisper (Speech-to-Text)}

Mari kita transkrip audio meeting menggunakan Python.

\begin{codeblock}[language=Python]
from openai import OpenAI
client = OpenAI()

audio_file = open("rekaman_rapat.mp3", "rb")
transcript = client.audio.transcriptions.create(
  model="whisper-1",
  file=audio_file
)

print(transcript.text)
\end{codeblock}
